% SPDX-FileCopyrightText: 2026 Yiyu Zhou <yiyu@yiyuzhou.io>
%
% SPDX-License-Identifier: GPL-3.0-or-later

\chapter{Keys}
\label{ch:keys}

\section{Keymaps}

cinq ships two keymaps.  The vi keymap, the default, moves with letters such
as \key{j} and \key{k}, and uses two-key sequences for some actions, such as
\key{d d}.  The cua keymap follows IBM's Common User Access, like most
desktop programs: \key{Control} with a letter edits, and \key{F1} opens the
help page.  Both put GO on \key{Space} and Panic on \key{Escape}, as QLab
does, and share the other keys used in this tutorial: \key{ArrowUp},
\key{ArrowDown}, \key{Enter}, \key{Alt} with those arrows, \key{Control+s},
\key{Control+o}, and \key{F11}.

The actions that this tutorial performed with buttons and clicks also have
keys, which differ between the keymaps:

\begin{center}
\begin{tabular}{lll}
\hline
Action & vi & cua \\
\hline
Add a cue below the selection & \key{o} & \key{Control+n} \\
Delete the selected cue & \key{d d} & \key{Delete} \\
Edit the name of the selected cue & \key{i} & \key{F2} \\
Undo the latest change & \key{u} & \key{Control+z} \\
Redo the change undone last & \key{Control+r} & \key{Control+y} \\
Show the next cue list, then the outputs & \key{l} & \key{Control+PageDown} \\
Show the next inspector tab & \key{g t} & \key{Control+Tab} \\
Show or hide the help & \key{?} & \key{F1} \\
\hline
\end{tabular}
\end{center}

\section{The help page}

\begin{enumerate}
\item Press the help key that the status bar names, \key{?} in the vi
  keymap.  The help page lists every action with its keys and its name in
  the configuration file, as Figure~\ref{fig:help} shows.
\item Press the same key again to close the page.  GO and Panic also close
  it, so that the show stays in view while it runs.
\end{enumerate}

\begin{figure}[H]
  \centering
  \includegraphics[width=\textwidth]{screenshots/help.png}
  \caption{Help page of the vi keymap}
  \label{fig:help}
\end{figure}

\section{Key sequences}

Some actions in the vi keymap take more than one key, such as \key{g g},
which selects the first cue.

\begin{enumerate}
\item Press \key{g}.  The status bar shows the keys typed so far, and a box
  lists the keys that can follow, as in Figure~\ref{fig:keys}.
\item Press \key{g} again.  The first cue is selected.
\end{enumerate}

\begin{figure}[H]
  \centering
  \includegraphics[width=\textwidth]{screenshots/keys.png}
  \caption{Keys that can follow \key{g}}
  \label{fig:keys}
\end{figure}

A number typed first repeats the action: \key{3 j} moves the selection three
rows down, and \key{3 d d} deletes three cues.  The keys for GO and Panic act
immediately, even in the middle of a sequence.

\section{Changing the keys}

The configuration file, \texttt{config.toml}, selects the keymap and binds
actions to other keys.  The help page links to it, and following the link
creates the file if it does not exist.  cinq rereads the file whenever it
changes, so edits take effect as soon as the file is saved.  The manual page
\texttt{cinq(5)} describes every setting and lists the keys of both keymaps.

\section*{Exercise}

You can now find the key for any action and choose the keys that suit you.
Switch to the cua keymap, as the examples in \texttt{cinq(5)} show, and open
the help page with \key{F1}.
